Une fente de largeur \qty{\v{larg}}{\micro\metre} est éclairée simultanément par deux lumières monochromatiques. L'une a une longueur d'onde de \qty{\v{lam1}}{nm}. On observe que le minimum de diffraction d'ordre \v{m1} de cette lumière coïncide avec le minimum d'ordre \v{m2} de l'autre lumière. L'écran est à \qty{\v{L}}{m} de la fente.

\begin{enumerate}[label=\alph*)]
    \item Déterminez la longueur d'onde de la seconde lumière.
    \item À quel angle se produit cette coïncidence ?
    \item À quelle distance du centre de l'écran se produit-elle ?
\end{enumerate}
